\documentclass{article}

\usepackage[dblblindworkshop]{neurips_2026}
\workshoptitle{Child Safety in AI}
\usepackage[utf8]{inputenc}
\usepackage[T1]{fontenc}
\usepackage{hyperref}
\usepackage{url}
\usepackage{booktabs}
\usepackage{amsmath}
\usepackage{amssymb}
\usepackage{xcolor}
\usepackage{microtype}

\title{Who Is Safe to Involve? ChildEsc Audits Escalation and Handoff Routing in Child-Facing AI}

\author{Anonymous Author(s)}

\begin{document}
\maketitle

\begin{abstract}
Child-facing AI must decide not only whether to escalate, but \emph{who is safe to involve}. Action-only evaluation can therefore reward a route that selects the right urgency while directing a child toward an inappropriate recipient. We introduce ChildEsc, a works-in-progress testbed of 80 LLM-assisted synthetic conversations in 20 counterfactual families. Each item specifies one of four actions plus permitted and forbidden handoff targets. We define the \emph{action-route gap}: exact action accuracy minus exact action-plus-target routing. A policy-exposed deterministic checksum achieves 65.0\% action accuracy but 53.8\% exact routing, an 11.3\% gap; 9/37 action-correct handoffs violate target constraints. A frozen relational contract audit passes 36/39 assertions across 7/9 families and retains two actionable failures. We also provide a provider-neutral, cache-replayable protocol for prompted-router audits, but report no live-model comparison because frozen credentials were unavailable. ChildEsc contributes an auditable routing specification and evidence gates, not a validated benchmark: practitioner and youth validation have not occurred, and no result supports deployment.
\end{abstract}

\section{Introduction}

Children use conversational AI for information, companionship, and emotional support, where a superficially safe response can still worsen a vulnerable situation \citep{unicef2026friend,cha2026beyond}. Binary respond/refuse policies are ill-suited to these interactions: refusal can close a help-seeking window, while unbounded engagement can reinforce dependency or miss imminent danger. Interviews with 19 youth-facing practitioners found that useful systems should gather context and bridge youth to tailored human support, but should not assume a parent is safe \citep{cha2026beyond}. The relevant safeguard decision is therefore \emph{what support action is proportionate now, and who is safe to involve?}

Related benchmarks cover important parts of this problem. MinorBench and Safe-Child-LLM emphasize unsafe-request handling \citep{khoo2025minorbench,jiao2025safechild}; KIDBench evaluates developmental response quality and trusted-adult redirection \citep{arif2026kidbench}; and CAREBench includes crisis referral, de-escalation, and human-support behavior across child-safety risks \citep{krishnakumar2026carebench}. MindEval measures multi-turn adult mental-health support \citep{pombal2025mindeval}. Most recently, CARE-Bench operationalized patient-facing medical triage as a four-label sequential current-action task, including whether more information is needed \citep{hua2026caretriage}. Thus neither ordered routing nor action timing is our novelty. ChildEsc instead isolates child-specific severity-matched trajectories and machine-checkable permitted and forbidden handoff sets. This focus maps directly to the workshop's restricted-data evaluation and deployment-safeguard themes \citep{childsafetyworkshop2026}.

We contribute: (1) a proposed four-action, context-constrained handoff specification instantiated as 80 synthetic trajectories; (2) metrics and item-level audits that expose action-correct but target-invalid routes; and (3) a reproducible checksum, provider-neutral prompted-router harness, source-to-design register, and concrete validation protocol. We deliberately do not claim benchmark validity, child-centeredness, or system superiority before independent review.

\section{ChildEsc testbed}

\paragraph{Scope.} ChildEsc applies to text-based consumer assistants, companions, tutors, and wellness or mental-health-adjacent systems used by or accessible to people under 18. It evaluates a router before or alongside response generation. It does not diagnose, recommend treatment, assess counseling language, or automatically contact any person or service.

\paragraph{Construction and provenance.} Under a written specification, an LLM coding assistant drafted 20 families across self-harm, acute distress, eating-related medical risk, abuse/coercion, online exploitation, AI dependency, unsafe requests, and privacy coercion. Each family has a shared context and four severity cues, expanded deterministically into three-user-turn conversations. No text derives from a real child, hotline transcript, clinical record, or explicit abuse material. A dated register maps domain-level design rationales to practitioner findings \citep{cha2026beyond}, child-rights guidance \citep{unicef2025guidance,unicef2026friend}, and crisis principles of collaboration and least-invasive intervention \citep{lifeline2024policy}. This retrospective mapping is \emph{source-informed}, not item-level grounding. The same workflow drafted labels and rules, so results are policy-exposed and unsuitable for out-of-sample claims.

\paragraph{Actions and targets.} The ordinal action space is $A=\{$\texttt{continue\_support} (0), \texttt{boundary\_redirect} (1), \texttt{human\_support} (2), \texttt{urgent\_handoff} (3)$\}$. Escalated items additionally permit one or more target classes: safe adult, non-implicated safe adult, school support, health professional, crisis service, emergency service, child protection, or platform reporting. Items can forbid an emittable target; for example, \texttt{safe\_adult} is forbidden when the implicated adult appears in the conversation, while \texttt{safe\_adult\_not\_implicated} may be permitted.

\paragraph{Failure-oriented metrics.} Let $g_i,p_i\in\{0,1,2,3\}$ be reference and predicted actions and $E=\{i:g_i\ge2\}$. Under-escalation is $|E|^{-1}\sum_{i\in E}\mathbb{1}[p_i<g_i]$. Urgent recall separately counts $3\rightarrow3$ decisions. A target is valid only when the predicted set is nonempty, is a subset of permitted targets, and excludes forbidden targets; targetless actions must emit none. Exact routing requires both the exact action and a valid target decision. The action-route gap is action accuracy minus exact-route accuracy. We also report counterfactual sensitivity and monotonicity across adjacent family variants. Reaction delay is omitted because v0.1 fixes the decisive cue turn.

\section{Diagnostic implementation and audit}

\paragraph{Policies and evaluator.} \textsc{Keyword} reacts only to explicit self-harm terms. \textsc{Agnostic} routes detected distress to non-urgent human support and dependency, exploitation, or unsafe-request cues to a boundary. \textsc{ChildEsc-Rules} uses phrase evidence, guards, conversation accumulation, domain rules, and conversation-derived adult-implication cues. A provider-neutral harness sends only conversation text and requests the same action-target schema from Gemini, OpenAI, Anthropic, or OpenRouter. It supports independent trial IDs, append-only attempt logs, strict parsing, and cache-only replay; invalid outputs make a run incomplete. No live provider results are reported because credentials were unavailable at protocol freeze. The checksum and prompted evaluator generate no counseling text, so supportive-response quality is outside scope.

\begin{table}[t]
  \caption{Policy-exposed diagnostics on 80 synthetic items. Action and Exact use all 80; Under and Valid use the 58 escalation-required items. Gap is Action minus Exact. These are implementation checks, not comparative effectiveness estimates.}
  \label{tab:results}
  \centering
  \small
  \begin{tabular}{lrrrrrr}
    \toprule
    Policy & Action & Urgent & Under & Valid & Gap & Exact \\
    \midrule
    Keyword & .225 & .189 & .879 & .121 & .000 & .225 \\
    Agnostic & .250 & .000 & .862 & .345 & .013 & .238 \\
    ChildEsc-Rules & .650 & .757 & .328 & .552 & .113 & .538 \\
    \bottomrule
  \end{tabular}
\end{table}

\paragraph{The action-route gap.} This diagnostic checksum verifies that generation, routing, constraints, and metrics execute end to end; it does not establish policy superiority. ChildEsc-Rules produces the exact action on 52/80 items but the exact action-plus-target route on only 43/80. Thus action-only scoring hides nine failures, an 11.3-point action-route gap. Conditioned on its 37 action-correct handoffs, 9/37 (24.3\%) violate target constraints. Family-bootstrap intervals are [.537, .762] for action accuracy, [.062, .163] for the gap, and [.438, .638] for exact routing. These intervals describe variation across the 20 selected families, not population uncertainty or protection from policy tuning.

\paragraph{Concrete failure.} In an abuse/coercion family, the reference route permits a non-implicated adult and child protection while forbidding the generic \texttt{safe\_adult}; a generic-recipient prediction can match \texttt{human\_support} yet fail the route. Across all escalation-required items, checksum targets are valid for 32, missing for 17, and contain extra unpermitted targets for nine. Two targetless items receive spurious targets. The benchmark therefore makes the recipient choice inspectable rather than treating ``escalate'' as sufficient.

\paragraph{Developer-specified stress tests.} These tests probe the frozen checksum without retuning and are not held-out evidence. Using only the latest user turn under-escalates all 58 escalation-required items. Disabling adult-implication inference leaves actions unchanged but emits forbidden \texttt{safe\_adult} in 7/12 implicated-adult cases. Uppercasing preserves every route; doubled spaces reduce action agreement to 45.0\% and urgent recall from 75.7\% to 10.8\%. Full results appear in Appendix~\ref{app:diagnostics}.

\section{Source-informed relational contract audit}

Aggregate labels cannot state how a route should change when one operational factor changes. We therefore froze nine matched, non-procedural probe families after v0.1 and before their first execution. Public sources motivate each relation---for example, temporal immediacy and least-invasive collaboration \citep{lifeline2024policy}, non-implicated support \citep{cha2026beyond}, and relational dependency \citep{unicef2026friend}---but do not supply ground-truth labels. Each family asserts an action relation plus atomic target, domain, or response-requirement constraints. The author had inspected the router; this is a post-freeze mechanistic audit, \emph{not held-out} evaluation.

\nolinenumbers
\begin{center}
  \refstepcounter{table}\label{tab:contracts}
  \parbox{.95\linewidth}{\small Table~\thetable: Frozen contract audit. Arrows show observed base-to-variant actions; Pass is atomic assertions. Family pass requires all assertions.}
  \vspace{4pt}
  {\scriptsize
  \begin{tabular}{@{}llll@{}}
    \toprule
    Controlled factor & Observed actions & Pass & Family \\
    \midrule
    Temporal immediacy & Urgent$\rightarrow$Urgent & 3/5 & fail \\
    Resolved negation & Urgent$\rightarrow$Continue & 3/3 & pass \\
    First/third person & Urgent$\rightarrow$Human & 4/4 & pass \\
    Implicated adult & Human$\rightarrow$Human & 6/6 & pass \\
    AI dependency & Continue$\rightarrow$Human & 4/4 & pass \\
    Coercive secrecy & Continue$\rightarrow$Human & 4/4 & pass \\
    Online threat & Boundary$\rightarrow$Urgent & 5/5 & pass \\
    Medical deterioration & Boundary$\rightarrow$Urgent & 4/5 & fail \\
    Explicit/implicit age & Human$\rightarrow$Human & 3/3 & pass \\
    \bottomrule
  \end{tabular}
  }
\end{center}
\linenumbers

Overall, 36/39 assertions and 7/9 families pass. The two failures are actionable and distinct from benchmark accuracy: the phrase ``safe right now'' triggers urgency despite explicit negation because \texttt{right now} is matched globally; physical deterioration triggers urgency but omits \texttt{health\_professional}. We retain both failures and do not modify the checksum. This illustrates ChildEsc's presently supported use: auditing whether an implementation conforms to an inspectable routing contract, not ranking models or certifying safety.

\paragraph{Pre-registered prompted-router audit.} Before inspecting any provider output, we froze exact candidate model IDs, prompt and schema hashes, three independent trials per model, a 256-token output cap, and omission of temperature. Each item is one API attempt: hidden retries and output repair are disallowed, and any provider error, refusal, or invalid schema makes the trial incomplete rather than silently shrinking its denominator. Complete trials are summarized by trial means and ranges; action and full-route instability are counted per item. Model differences use candidate-minus-reference family-paired bootstrap intervals, not independent-row tests. Content-addressed caches include the visible request, trial ID, and response hashes but exclude credentials; a cache-only replay must account for every item. The harness has deterministic transport tests for Gemini, OpenAI, Anthropic, and OpenRouter. Because no credential was available at freeze, we invoke the pre-specified checksum fallback and make no prompted-model performance claim.

\section{Validity, impact, and next evaluation}

The central validity threat is circularity: scenarios, initial labels, lexical rules, probes, and prose were drafted in one LLM-assisted workflow. Source traceability improves auditability, not validity. The small English convenience set omits dialect, code-switching, disability, culture, long relationships, and adversarial obfuscation. Severity is operational rather than clinical; v0.1 also does not isolate information-seeking as an action, unlike CARE-Bench. Target validity measures taxonomy conformance, not accessible, consensual, locally available, or successful handoff. No practitioner or youth validation has occurred. ``Child-centered'' is therefore a design objective, not a measured property.

Before comparative or generalization claims, qualified practitioners must independently review the construct, action thresholds, target constraints, missing context, and jurisdictional feasibility; youth advisors must separately review plausibility, agency, and accessibility without personal disclosure. Disagreement and pre-review labels will be retained. Only then will independent authors create family-locked public and controlled-access tests without seeing the router. This sequence follows benchmark-validity guidance to connect phenomenon, task, items, metrics, and claims while separating development from test evidence \citep{bean2025construct}.

\paragraph{Evidence gates.} Gate 1 versions practitioner revisions rather than overwriting v0.1 and reports label distributions, uncertainty, and disagreement. Gate 2 keeps youth judgments of agency and plausibility separate from professional risk judgments. Gate 3 locks duplicate clusters by family, reserves a controlled-access test, and evaluates it once without router tuning. Gate 4 audits any mapper from open-ended responses into ChildEsc actions against independent human coding. Failure at any gate blocks claims that depend on the next one; no gate alone certifies deployment safety.

\paragraph{Prospective validation outputs.} Practitioner review would produce independent action and permitted/forbidden-target labels, missing-context flags, and jurisdictional-feasibility notes; youth review would separately produce plausibility, agency, accessibility, and handoff-acceptability judgments without personal disclosure. ChildEsc would report pre-review labels, revisions, uncertainty, and disagreement rather than collapsing all perspectives into a consensus gold label. Controlled-access families would be authored by people unexposed to router rules and evaluated once. These steps are a pre-specified protocol, not completed validation.

ChildEsc's practical use is deliberately narrow: give developers and child-safety reviewers an executable object for debating escalation thresholds, unsafe-default recipients, and what evidence is missing before a handoff. The current artifact can audit schema conformance and surface disagreements; it cannot determine clinical risk, select a real person, assess response quality, or authorize disclosure.

\paragraph{Questions for the workshop.} (1) When context is insufficient, should a router clarify, abstain, or emit multiple consent-preserving routes rather than force one label? (2) Which judgments belong separately to practitioners (risk and service feasibility) and youth advisors (plausibility, agency, and acceptability), and how should disagreement appear in scoring? (3) What evidence and governance must precede any operational handoff, especially when a nominally trusted adult may be implicated? Answers would directly determine ChildEsc v0.2's action space, annotation protocol, and release gates.

\clearpage
\bibliographystyle{plainnat}
\bibliography{references}

\appendix

\section{Artifact and schema}

The anonymous supplement contains standard-library Python source, 48 tests for the frozen scientific and manuscript path, 22 provider-path tests, source families, a dated source register, retrospective family rationale map, controlled authoring protocol, frozen contract probes, generated data, item-level predictions and assertions, aggregate metrics, error analyses, ablations, robustness checks, a gated follow-up specification, and exact paper tables. Running \texttt{make reproduce} regenerates all 80 items and every reported value. Regression tests verify input boundaries, source IDs, complete family mapping, safe follow-up defaults, and synchronization between versioned results and the manuscript.

\section{Relationship to closest benchmarks}

CAREBench already evaluates crisis referral, de-escalation, and redirection toward human support; KIDBench evaluates multi-turn developmental quality, boundary-setting, and trusted-adult redirection. CARE-Bench already evaluates four ordered current-action labels over controlled patient-disclosure prefixes and demonstrates the importance of an information-needed state. ChildEsc does not claim these behaviors, ordered actions, or timing as new. Its narrower proposed unit is a child-specific route over severity-matched counterfactuals coupled to set-valued permitted and forbidden handoff targets. Whether that decomposition is useful or child-centered remains an empirical question.

\section{Representative examples}

\begin{table}[h]
  \caption{Abbreviated synthetic examples; no procedural harmful detail is included.}
  \centering
  \small
  \begin{tabular}{@{}p{.13\linewidth}p{.41\linewidth}p{.18\linewidth}p{.18\linewidth}@{}}
    \toprule
    Domain & Child cue & Reference action & Permitted target example \\
    \midrule
    Dependency & ``Can you be my best friend forever?'' & Boundary & none \\
    Self-harm & ``I don't care if I wake up tomorrow.'' & Human support & safe adult or professional \\
    Abuse & ``The adult is outside my locked door now.'' & Urgent & non-implicated adult \\
    Eating health & ``I fainted and my chest feels strange.'' & Urgent & health or emergency service \\
    \bottomrule
  \end{tabular}
\end{table}

\section{Metric details}

Unsafe-target rate is conditioned on items with at least one forbidden target. Strict valid-target rate is conditioned on items with at least one permitted target. Exact route requires an exact action match; if permitted targets exist, the predicted set must be nonempty and a subset of the permitted set, with no forbidden intersection; otherwise it must be empty. The action-route gap equals exact action accuracy minus exact-route accuracy; the companion failure rate conditions target-invalid routes on action-correct handoffs. Counterfactual sensitivity is conditioned on adjacent family transitions where the reference action increases. Family bootstrap resampling keeps each four-item family intact and uses seed 20260829 for 1,000 repetitions.

\section{Diagnostic failure analysis}
\label{app:diagnostics}

Table~\ref{tab:diagnostics} reports every developer-specified input-evidence ablation and orthographic transformation. Agree is action agreement with the untransformed full-conversation run. Under is conditioned on the 58 escalation-required items, Valid on those with permitted targets, and Unsafe on the 12 items with a forbidden target. The latest-turn condition is a deliberately destructive evidence ablation because every decisive cue occurs on the second user turn. Adult inference is disabled in a scoped call and restored afterward, leaving action logic and conversation text unchanged. Orthographic tests are developer-authored, policy-exposed diagnostics rather than paraphrase or dialect generalization tests.

\begin{table}[h]
  \caption{Frozen-router diagnostic analyses.}
  \label{tab:diagnostics}
  \centering
  \scriptsize
  \begin{tabular}{lrrrrrr}
    \toprule
    Condition & Agree & Urgent & Under & Valid & Unsafe & Exact \\
    \midrule
    Full conversation & 1.000 & .757 & .328 & .552 & .000 & .538 \\
    Latest turn only & .338 & .000 & 1.000 & .000 & .000 & .138 \\
    Adult inference off & 1.000 & .757 & .328 & .431 & .583 & .488 \\
    Uppercase & 1.000 & .757 & .328 & .552 & .000 & .538 \\
    Extra spaces & .450 & .108 & .897 & .034 & .083 & .175 \\
    Expanded contractions & .975 & .730 & .362 & .534 & .000 & .525 \\
    \bottomrule
  \end{tabular}
\end{table}

The item-level taxonomy records 240 policy-item rows. Three original guard probes verify coded invariants. The post-freeze contract suite adds nine matched families and 39 atomic assertions spanning action order, targets, domains, and response requirements. Its JSON source, item-level CSV, and aggregate table retain all failures. Passing authored probes is not a performance estimate.

\section{Source traceability and future authoring}

The source register records what each source supports and explicitly does not support. The v0.1 family map is retrospective and domain-level; it does not make individual cases source-grounded. The prospective authoring protocol requires one-factor isolation, information-state timing, age and locale metadata, source-to-case transformation rationale, non-procedural content, independent first-pass labels without router output, retention of disagreement, and family-locked development, validation, public-test, and controlled-access splits. Items with missing context may remain unresolved rather than receiving a forced label.

\section{Prospective validation protocol}

No external ethics determination or approved human-participant protocol is documented for the reported work; therefore no recruitment or data collection occurred. The supplement provides prospective practitioner and youth materials, a data-management plan, and a safeguarding template. Subject to appropriate approval, practitioners would independently label action, permitted and forbidden targets, missing context, and coercion risk using a balanced incomplete-block design. Youth advisors would review plausibility, agency, accessibility, and handoff framing without being asked to disclose personal experiences. Disagreement and pre-review labels would be retained. Protocol preparation is not validation.

\section{Prospective delegated follow-up extension}

The supplement formalizes the proposed 1--24 hour crisis-service follow-up as a non-operational governance specification. It is not an established 988 policy, and no cutoff is validated. It requires a written 988/center partnership, independent safeguarding and legal review, separate revocable consent that does not condition account access, human crisis-professional approval, minimum-necessary transfer, receiving-service capacity guarantees, and prospective measurement of false alerts, missed urgency, chilling effects, coercion, and disparities. The deadline function remains unspecified: a multidisciplinary committee including 988 operations, child-safety, clinical, privacy, disability, cultural, and youth representatives must pre-register and validate it. Imminent danger is excluded from a delayed queue. Until those gates pass, ChildEsc must not collect PII, schedule follow-up, transmit data, or initiate contact.

\section{Independent holdout status}

No independently authored holdout is reported. The supplement specifies a sealed authoring procedure requiring a contributor who has not inspected router rules, existing scenarios, or item-level failures. The current authors and coding agents are ineligible under that procedure. Any future holdout must be evaluated once without router tuning and reported separately from v0.1.

\section{Compute, ethics, and LLM usage}

Reported experiments use no GPU, model API, training, or external service and complete in under one second on a laptop-class CPU. A general-purpose LLM coding assistant drafted candidate scenarios and initial labels, implementation, literature notes, and manuscript text. No LLM judged outputs or generated reported measurements. The optional prompted-router harness was tested with deterministic mock transports, but no live provider result is reported. Public release must retain content warnings and must not support automated data disclosure, emergency or caregiver contact, or claims of clinical validation.

\clearpage
\section*{NeurIPS Paper Checklist}

\begin{enumerate}

\item {\bf Claims}
    \item[] Question: Do the main claims made in the abstract and introduction accurately reflect the paper's contributions and scope?
    \item[] Answer: \answerYes{}
    \item[] Justification: The abstract and Introduction frame ChildEsc as a preliminary testbed and specification, label the evaluation policy-exposed, avoid comparative-effectiveness claims, and state that validation is required.
    \item[] Guidelines:
    \begin{itemize}
        \item The answer \answerNA{} means that the abstract and introduction do not include the claims made in the paper.
        \item The abstract and/or introduction should clearly state the claims made, including the contributions made in the paper and important assumptions and limitations. A \answerNo{} or \answerNA{} answer to this question will not be perceived well by the reviewers. 
        \item The claims made should match theoretical and experimental results, and reflect how much the results can be expected to generalize to other settings. 
        \item It is fine to include aspirational goals as motivation as long as it is clear that these goals are not attained by the paper. 
    \end{itemize}

\item {\bf Limitations}
    \item[] Question: Does the paper discuss the limitations of the work performed by the authors?
    \item[] Answer: \answerYes{}
    \item[] Justification: The dedicated Validity section discusses policy exposure, lexical circularity, synthetic English data, construct and cultural validity, target abstraction, bootstrap scope, and absent human validation.
    \item[] Guidelines:
    \begin{itemize}
        \item The answer \answerNA{} means that the paper has no limitation while the answer \answerNo{} means that the paper has limitations, but those are not discussed in the paper. 
        \item The authors are encouraged to create a separate ``Limitations'' section in their paper.
        \item The paper should point out any strong assumptions and how robust the results are to violations of these assumptions (e.g., independence assumptions, noiseless settings, model well-specification, asymptotic approximations only holding locally). The authors should reflect on how these assumptions might be violated in practice and what the implications would be.
        \item The authors should reflect on the scope of the claims made, e.g., if the approach was only tested on a few datasets or with a few runs. In general, empirical results often depend on implicit assumptions, which should be articulated.
        \item The authors should reflect on the factors that influence the performance of the approach. For example, a facial recognition algorithm may perform poorly when image resolution is low or images are taken in low lighting. Or a speech-to-text system might not be used reliably to provide closed captions for online lectures because it fails to handle technical jargon.
        \item The authors should discuss the computational efficiency of the proposed algorithms and how they scale with dataset size.
        \item If applicable, the authors should discuss possible limitations of their approach to address problems of privacy and fairness.
        \item While the authors might fear that complete honesty about limitations might be used by reviewers as grounds for rejection, a worse outcome might be that reviewers discover limitations that aren't acknowledged in the paper. The authors should use their best judgment and recognize that individual actions in favor of transparency play an important role in developing norms that preserve the integrity of the community. Reviewers will be specifically instructed to not penalize honesty concerning limitations.
    \end{itemize}

\item {\bf Theory assumptions and proofs}
    \item[] Question: For each theoretical result, does the paper provide the full set of assumptions and a complete (and correct) proof?
    \item[] Answer: \answerNA{}
    \item[] Justification: The paper makes no theoretical claims and contains no theorem or proof.
    \item[] Guidelines:
    \begin{itemize}
        \item The answer \answerNA{} means that the paper does not include theoretical results. 
        \item All the theorems, formulas, and proofs in the paper should be numbered and cross-referenced.
        \item All assumptions should be clearly stated or referenced in the statement of any theorems.
        \item The proofs can either appear in the main paper or the supplemental material, but if they appear in the supplemental material, the authors are encouraged to provide a short proof sketch to provide intuition. 
        \item Inversely, any informal proof provided in the core of the paper should be complemented by formal proofs provided in appendix or supplemental material.
        \item Theorems and Lemmas that the proof relies upon should be properly referenced. 
    \end{itemize}

    \item {\bf Experimental result reproducibility}
    \item[] Question: Does the paper fully disclose all the information needed to reproduce the main experimental results of the paper to the extent that it affects the main claims and/or conclusions of the paper (regardless of whether the code and data are provided or not)?
    \item[] Answer: \answerYes{}
    \item[] Justification: The main paper and Appendix specify the data construction, policies, metrics, seed, bootstrap repetitions, and one-command reproduction path.
    \item[] Guidelines:
    \begin{itemize}
        \item The answer \answerNA{} means that the paper does not include experiments.
        \item If the paper includes experiments, a \answerNo{} answer to this question will not be perceived well by the reviewers: Making the paper reproducible is important, regardless of whether the code and data are provided or not.
        \item If the contribution is a dataset and\slash or model, the authors should describe the steps taken to make their results reproducible or verifiable. 
        \item Depending on the contribution, reproducibility can be accomplished in various ways. For example, if the contribution is a novel architecture, describing the architecture fully might suffice, or if the contribution is a specific model and empirical evaluation, it may be necessary to either make it possible for others to replicate the model with the same dataset, or provide access to the model. In general. releasing code and data is often one good way to accomplish this, but reproducibility can also be provided via detailed instructions for how to replicate the results, access to a hosted model (e.g., in the case of a large language model), releasing of a model checkpoint, or other means that are appropriate to the research performed.
        \item While NeurIPS does not require releasing code, the conference does require all submissions to provide some reasonable avenue for reproducibility, which may depend on the nature of the contribution. For example
        \begin{enumerate}
            \item If the contribution is primarily a new algorithm, the paper should make it clear how to reproduce that algorithm.
            \item If the contribution is primarily a new model architecture, the paper should describe the architecture clearly and fully.
            \item If the contribution is a new model (e.g., a large language model), then there should either be a way to access this model for reproducing the results or a way to reproduce the model (e.g., with an open-source dataset or instructions for how to construct the dataset).
            \item We recognize that reproducibility may be tricky in some cases, in which case authors are welcome to describe the particular way they provide for reproducibility. In the case of closed-source models, it may be that access to the model is limited in some way (e.g., to registered users), but it should be possible for other researchers to have some path to reproducing or verifying the results.
        \end{enumerate}
    \end{itemize}


\item {\bf Open access to data and code}
    \item[] Question: Does the paper provide open access to the data and code, with sufficient instructions to faithfully reproduce the main experimental results, as described in supplemental material?
    \item[] Answer: \answerYes{}
    \item[] Justification: The anonymous supplement includes source families, the source register, contract probes and assertions, generated data, standard-library code, tests, item predictions, and exact commands; public release remains subject to the documented safety review.
    \item[] Guidelines:
    \begin{itemize}
        \item The answer \answerNA{} means that paper does not include experiments requiring code.
        \item Please see the NeurIPS code and data submission guidelines (\url{https://neurips.cc/public/guides/CodeSubmissionPolicy}) for more details.
        \item While we encourage the release of code and data, we understand that this might not be possible, so \answerNo{} is an acceptable answer. Papers cannot be rejected simply for not including code, unless this is central to the contribution (e.g., for a new open-source benchmark).
        \item The instructions should contain the exact command and environment needed to run to reproduce the results. See the NeurIPS code and data submission guidelines (\url{https://neurips.cc/public/guides/CodeSubmissionPolicy}) for more details.
        \item The authors should provide instructions on data access and preparation, including how to access the raw data, preprocessed data, intermediate data, and generated data, etc.
        \item The authors should provide scripts to reproduce all experimental results for the new proposed method and baselines. If only a subset of experiments are reproducible, they should state which ones are omitted from the script and why.
        \item At submission time, to preserve anonymity, the authors should release anonymized versions (if applicable).
        \item Providing as much information as possible in supplemental material (appended to the paper) is recommended, but including URLs to data and code is permitted.
    \end{itemize}


\item {\bf Experimental setting/details}
    \item[] Question: Does the paper specify all the training and test details (e.g., data splits, hyperparameters, how they were chosen, type of optimizer) necessary to understand the results?
    \item[] Answer: \answerYes{}
    \item[] Justification: Section 3 and the Appendix specify the three deterministic policies, the absence of training or tuning splits, and the fact that policy exposure prevents out-of-sample interpretation.
    \item[] Guidelines:
    \begin{itemize}
        \item The answer \answerNA{} means that the paper does not include experiments.
        \item The experimental setting should be presented in the core of the paper to a level of detail that is necessary to appreciate the results and make sense of them.
        \item The full details can be provided either with the code, in appendix, or as supplemental material.
    \end{itemize}

\item {\bf Experiment statistical significance}
    \item[] Question: Does the paper report error bars suitably and correctly defined or other appropriate information about the statistical significance of the experiments?
    \item[] Answer: \answerYes{}
    \item[] Justification: The audit reports 95\% percentile intervals from 1,000 seeded family-level bootstrap resamples and explicitly limits them to variation across the 20 selected families, not population uncertainty or tuning bias.
    \item[] Guidelines:
    \begin{itemize}
        \item The answer \answerNA{} means that the paper does not include experiments.
        \item The authors should answer \answerYes{} if the results are accompanied by error bars, confidence intervals, or statistical significance tests, at least for the experiments that support the main claims of the paper.
        \item The factors of variability that the error bars are capturing should be clearly stated (for example, train/test split, initialization, random drawing of some parameter, or overall run with given experimental conditions).
        \item The method for calculating the error bars should be explained (closed form formula, call to a library function, bootstrap, etc.)
        \item The assumptions made should be given (e.g., Normally distributed errors).
        \item It should be clear whether the error bar is the standard deviation or the standard error of the mean.
        \item It is OK to report 1-sigma error bars, but one should state it. The authors should preferably report a 2-sigma error bar than state that they have a 96\% CI, if the hypothesis of Normality of errors is not verified.
        \item For asymmetric distributions, the authors should be careful not to show in tables or figures symmetric error bars that would yield results that are out of range (e.g., negative error rates).
        \item If error bars are reported in tables or plots, the authors should explain in the text how they were calculated and reference the corresponding figures or tables in the text.
    \end{itemize}

\item {\bf Experiments compute resources}
    \item[] Question: For each experiment, does the paper provide sufficient information on the computer resources (type of compute workers, memory, time of execution) needed to reproduce the experiments?
    \item[] Answer: \answerYes{}
    \item[] Justification: The Compute appendix states that evaluation uses a laptop-class CPU, no GPU or external API, and completes in under one second in the reported environment.
    \item[] Guidelines:
    \begin{itemize}
        \item The answer \answerNA{} means that the paper does not include experiments.
        \item The paper should indicate the type of compute workers CPU or GPU, internal cluster, or cloud provider, including relevant memory and storage.
        \item The paper should provide the amount of compute required for each of the individual experimental runs as well as estimate the total compute. 
        \item The paper should disclose whether the full research project required more compute than the experiments reported in the paper (e.g., preliminary or failed experiments that didn't make it into the paper). 
    \end{itemize}
    
\item {\bf Code of ethics}
    \item[] Question: Does the research conducted in the paper conform, in every respect, with the NeurIPS Code of Ethics \url{https://neurips.cc/public/EthicsGuidelines}?
    \item[] Answer: \answerYes{}
    \item[] Justification: The work uses only synthetic, non-procedural text and no human subjects; sensitive-content, release, dual-use, and overclaiming risks are disclosed.
    \item[] Guidelines:
    \begin{itemize}
        \item The answer \answerNA{} means that the authors have not reviewed the NeurIPS Code of Ethics.
        \item If the authors answer \answerNo, they should explain the special circumstances that require a deviation from the Code of Ethics.
        \item The authors should make sure to preserve anonymity (e.g., if there is a special consideration due to laws or regulations in their jurisdiction).
    \end{itemize}


\item {\bf Broader impacts}
    \item[] Question: Does the paper discuss both potential positive societal impacts and negative societal impacts of the work performed?
    \item[] Answer: \answerYes{}
    \item[] Justification: Section 5 discusses privacy-preserving audit benefits and risks of surveillance, coercive intervention, cultural mismatch, false safety certification, and unsolicited crisis-data transfer, with explicit consent and evidence gates.
    \item[] Guidelines:
    \begin{itemize}
        \item The answer \answerNA{} means that there is no societal impact of the work performed.
        \item If the authors answer \answerNA{} or \answerNo, they should explain why their work has no societal impact or why the paper does not address societal impact.
        \item Examples of negative societal impacts include potential malicious or unintended uses (e.g., disinformation, generating fake profiles, surveillance), fairness considerations (e.g., deployment of technologies that could make decisions that unfairly impact specific groups), privacy considerations, and security considerations.
        \item The conference expects that many papers will be foundational research and not tied to particular applications, let alone deployments. However, if there is a direct path to any negative applications, the authors should point it out. For example, it is legitimate to point out that an improvement in the quality of generative models could be used to generate Deepfakes for disinformation. On the other hand, it is not needed to point out that a generic algorithm for optimizing neural networks could enable people to train models that generate Deepfakes faster.
        \item The authors should consider possible harms that could arise when the technology is being used as intended and functioning correctly, harms that could arise when the technology is being used as intended but gives incorrect results, and harms following from (intentional or unintentional) misuse of the technology.
        \item If there are negative societal impacts, the authors could also discuss possible mitigation strategies (e.g., gated release of models, providing defenses in addition to attacks, mechanisms for monitoring misuse, mechanisms to monitor how a system learns from feedback over time, improving the efficiency and accessibility of ML).
    \end{itemize}
    
\item {\bf Safeguards}
    \item[] Question: Does the paper describe safeguards that have been put in place for responsible release of data or models that have a high risk for misuse (e.g., pre-trained language models, image generators, or scraped datasets)?
    \item[] Answer: \answerYes{}
    \item[] Justification: The Dataset Card, Ethical release notes, and prospective validation protocols require content warnings, prohibit automated data disclosure, emergency or caregiver contact, and clinical claims, and label current outputs as unvalidated policy-exposed diagnostics.
    \item[] Guidelines:
    \begin{itemize}
        \item The answer \answerNA{} means that the paper poses no such risks.
        \item Released models that have a high risk for misuse or dual-use should be released with necessary safeguards to allow for controlled use of the model, for example by requiring that users adhere to usage guidelines or restrictions to access the model or implementing safety filters. 
        \item Datasets that have been scraped from the Internet could pose safety risks. The authors should describe how they avoided releasing unsafe images.
        \item We recognize that providing effective safeguards is challenging, and many papers do not require this, but we encourage authors to take this into account and make a best faith effort.
    \end{itemize}

\item {\bf Licenses for existing assets}
    \item[] Question: Are the creators or original owners of assets (e.g., code, data, models), used in the paper, properly credited and are the license and terms of use explicitly mentioned and properly respected?
    \item[] Answer: \answerNA{}
    \item[] Justification: No external dataset, model, or code asset is used in the experiments; prior scholarship and public guidance are cited but not redistributed.
    \item[] Guidelines:
    \begin{itemize}
        \item The answer \answerNA{} means that the paper does not use existing assets.
        \item The authors should cite the original paper that produced the code package or dataset.
        \item The authors should state which version of the asset is used and, if possible, include a URL.
        \item The name of the license (e.g., CC-BY 4.0) should be included for each asset.
        \item For scraped data from a particular source (e.g., website), the copyright and terms of service of that source should be provided.
        \item If assets are released, the license, copyright information, and terms of use in the package should be provided. For popular datasets, \url{paperswithcode.com/datasets} has curated licenses for some datasets. Their licensing guide can help determine the license of a dataset.
        \item For existing datasets that are re-packaged, both the original license and the license of the derived asset (if it has changed) should be provided.
        \item If this information is not available online, the authors are encouraged to reach out to the asset's creators.
    \end{itemize}

\item {\bf New assets}
    \item[] Question: Are new assets introduced in the paper well documented and is the documentation provided alongside the assets?
    \item[] Answer: \answerYes{}
    \item[] Justification: The supplement includes a README, specification, dataset card, schema, provenance, intended and excluded uses, limitations, tests, and reproduction commands.
    \item[] Guidelines:
    \begin{itemize}
        \item The answer \answerNA{} means that the paper does not release new assets.
        \item Researchers should communicate the details of the dataset\slash code\slash model as part of their submissions via structured templates. This includes details about training, license, limitations, etc. 
        \item The paper should discuss whether and how consent was obtained from people whose asset is used.
        \item At submission time, remember to anonymize your assets (if applicable). You can either create an anonymized URL or include an anonymized zip file.
    \end{itemize}

\item {\bf Crowdsourcing and research with human subjects}
    \item[] Question: For crowdsourcing experiments and research with human subjects, does the paper include the full text of instructions given to participants and screenshots, if applicable, as well as details about compensation (if any)? 
    \item[] Answer: \answerNA{}
    \item[] Justification: The reported work includes no crowdsourcing or human-subject research; the proposed future validation is clearly described as not yet conducted.
    \item[] Guidelines:
    \begin{itemize}
        \item The answer \answerNA{} means that the paper does not involve crowdsourcing nor research with human subjects.
        \item Including this information in the supplemental material is fine, but if the main contribution of the paper involves human subjects, then as much detail as possible should be included in the main paper. 
        \item According to the NeurIPS Code of Ethics, workers involved in data collection, curation, or other labor should be paid at least the minimum wage in the country of the data collector. 
    \end{itemize}

\item {\bf Institutional review board (IRB) approvals or equivalent for research with human subjects}
    \item[] Question: Does the paper describe potential risks incurred by study participants, whether such risks were disclosed to the subjects, and whether Institutional Review Board (IRB) approvals (or an equivalent approval/review based on the requirements of your country or institution) were obtained?
    \item[] Answer: \answerNA{}
    \item[] Justification: No human-subject research was conducted. Future practitioner and youth-participatory work will undergo appropriate institutional ethics review before recruitment.
    \item[] Guidelines:
    \begin{itemize}
        \item The answer \answerNA{} means that the paper does not involve crowdsourcing nor research with human subjects.
        \item Depending on the country in which research is conducted, IRB approval (or equivalent) may be required for any human subjects research. If you obtained IRB approval, you should clearly state this in the paper. 
        \item We recognize that the procedures for this may vary significantly between institutions and locations, and we expect authors to adhere to the NeurIPS Code of Ethics and the guidelines for their institution. 
        \item For initial submissions, do not include any information that would break anonymity (if applicable), such as the institution conducting the review.
    \end{itemize}

\item {\bf Declaration of LLM usage}
    \item[] Question: Does the paper describe the usage of LLMs if it is an important, original, or non-standard component of the core methods in this research? Note that if the LLM is used only for writing, editing, or formatting purposes and does \emph{not} impact the core methodology, scientific rigor, or originality of the research, declaration is not required.
    %this research? 
    \item[] Answer: \answerYes{}
    \item[] Justification: The Declaration of LLM usage states that an LLM assistant drafted candidate scenarios, initial labels, implementation, and prose; no LLM judged outputs or generated reported measurements.
    \item[] Guidelines:
    \begin{itemize}
        \item The answer \answerNA{} means that the core method development in this research does not involve LLMs as any important, original, or non-standard components.
        \item Please refer to our LLM policy in the NeurIPS handbook for what should or should not be described.
    \end{itemize}

\end{enumerate}


\end{document}
